% GENERATED FILE. Edit canonical lessons, then run npm run generate:books.
% canonical-source-hash: fnv1a64:d46943bb55c31899
% canonical-lessons: BN-C350-adorsho, BN-C350-mulyobodh, BN-C350-sohishnuta, BN-C350-sotyobadi, BN-C350-nejyo

\chapter{What We Value}
\label{ch:bn-what-we-value}
\hlchaptermodality{350}

\begin{chapteropening}
\textbf{By the end of this chapter:} \emph{I can say what I value and why.}

Complete the last lesson of chapter 350: I can say what I value and why.
\end{chapteropening}

\section[\texorpdfstring{\bn{আদর্শ} (ādôrshô)}{আদর্শ (ādôrshô)}]{\bn{আদর্শ} (ādôrshô) — an ideal; a model to follow}

\label{lesson:BN-C350-adorsho}



\begin{quote}
Before the new one: say the Bengali for a policy, then the Bengali for population.
\end{quote}

\subsection*{You'll want to know: \bn{আদর্শ}}
\textbf{\bn{আদর্শ}} — \emph{ādôrshô} — \textquotedblleft{}an ideal; a model to follow\textquotedblright{}.

\begin{quote}
\bn{সত্যি} \bn{বলা} \bn{আমাদের} \bn{পরিবারের} \bn{আদর্শ।}
\end{quote}

\begin{quote}
\emph{sôtti bôlā āmāder pôribārer ādôrshô.}
\end{quote}

\begin{quote}
Telling the truth is our family's ideal.
\end{quote}

\begin{cousinweb}
\textbf{\bn{আদর্শ}} (\emph{ādôrshô}) is Sanskrit \emph{ādarśa}, a mirror (Monier-Williams). The sense 'an ideal', the image you shape yourself by, is a later one.
\end{cousinweb}

\begin{grammarlens}[title={What you've built}]
One more word for giving reasons and opinions.
\end{grammarlens}

\subsection*{Your turn}
\begin{itemize}
\raggedright
  \item \emph{Say it:} \emph{ādôrshô}
  \item \emph{Say it:} \emph{ādôrshô}, once more
  \item \emph{Say it:} \emph{sôtti bôlā āmāder pôribārer ādôrshô.} — Telling the truth is our family's ideal.
  \item \emph{Recall:} say the Bengali for a policy, then the Bengali for population, then say \emph{ādôrshô} again
\end{itemize}

\begin{tcolorbox}[breakable,colback=teal!4,colframe=teal!35!black,title={Before you move on}]
What does \bn{আদর্শ} mean? (\textquotedblleft{}An ideal; a model to follow\textquotedblright{}.) Say it once more.
\end{tcolorbox}
\section[\texorpdfstring{\bn{মূল্যবোধ} (mūlyôbodh)}{মূল্যবোধ (mūlyôbodh)}]{\bn{মূল্যবোধ} (mūlyôbodh) — moral values}

\label{lesson:BN-C350-mulyobodh}



\begin{quote}
Before the new one: say the Bengali for population, then the Bengali for an ideal.
\end{quote}

\subsection*{You'll want to know: \bn{মূল্যবোধ}}
\textbf{\bn{মূল্যবোধ}} — \emph{mūlyôbodh} — \textquotedblleft{}moral values\textquotedblright{}.

\begin{quote}
\bn{পরিবার} \bn{থেকেই} \bn{আমরা} \bn{মূল্যবোধ} \bn{শিখি।}
\end{quote}

\begin{quote}
\emph{pôribār thekei āmrā mūlyôbodh shikhi.}
\end{quote}

\begin{quote}
We learn our values from our families.
\end{quote}

\begin{cousinweb}
\textbf{\bn{মূল্যবোধ}} (\emph{mūlyôbodh}) is a modern Bengali compound of Sanskrit \emph{mūlya}, value or price, and \emph{bodha}, understanding (Monier-Williams): a sense of what things are worth.
\end{cousinweb}

\begin{grammarlens}[title={What you've built}]
One more word for giving reasons and opinions.
\end{grammarlens}

\subsection*{Your turn}
\begin{itemize}
\raggedright
  \item \emph{Say it:} \emph{mūlyôbodh}
  \item \emph{Say it:} \emph{mūlyôbodh}, once more
  \item \emph{Say it:} \emph{pôribār thekei āmrā mūlyôbodh shikhi.} — We learn our values from our families.
  \item \emph{Recall:} say the Bengali for population, then the Bengali for an ideal, then say \emph{mūlyôbodh} again
\end{itemize}

\begin{tcolorbox}[breakable,colback=teal!4,colframe=teal!35!black,title={Before you move on}]
What does \bn{মূল্যবোধ} mean? (\textquotedblleft{}Moral values\textquotedblright{}.) Say it once more.
\end{tcolorbox}
\section[\texorpdfstring{\bn{সহিষ্ণুতা} (sôhishnutā)}{সহিষ্ণুতা (sôhishnutā)}]{\bn{সহিষ্ণুতা} (sôhishnutā) — tolerance}

\label{lesson:BN-C350-sohishnuta}



\begin{quote}
Before the new one: say the Bengali for an ideal, then the Bengali for moral values.
\end{quote}

\subsection*{You'll want to know: \bn{সহিষ্ণুতা}}
\textbf{\bn{সহিষ্ণুতা}} — \emph{sôhishnutā} — \textquotedblleft{}tolerance\textquotedblright{}.

\begin{quote}
\bn{সমাজে} \bn{সহিষ্ণুতা} \bn{দরকার।}
\end{quote}

\begin{quote}
\emph{sômāje sôhishnutā dôrkār.}
\end{quote}

\begin{quote}
Society needs tolerance.
\end{quote}

\begin{cousinweb}
\textbf{\bn{সহিষ্ণুতা}} (\emph{sôhishnutā}) is Sanskrit \emph{sahiṣṇutā}, the abstract of \emph{sahiṣṇu}, patient and forbearing (Monier-Williams): the patience to live with views you do not share.
\end{cousinweb}

\begin{grammarlens}[title={What you've built}]
One more word for giving reasons and opinions.
\end{grammarlens}

\subsection*{Your turn}
\begin{itemize}
\raggedright
  \item \emph{Say it:} \emph{sôhishnutā}
  \item \emph{Say it:} \emph{sôhishnutā}, once more
  \item \emph{Say it:} \emph{sômāje sôhishnutā dôrkār.} — Society needs tolerance.
  \item \emph{Recall:} say the Bengali for an ideal, then the Bengali for moral values, then say \emph{sôhishnutā} again
\end{itemize}

\begin{tcolorbox}[breakable,colback=teal!4,colframe=teal!35!black,title={Before you move on}]
What does \bn{সহিষ্ণুতা} mean? (\textquotedblleft{}Tolerance\textquotedblright{}.) Say it once more.
\end{tcolorbox}
\section[\texorpdfstring{\bn{সত্যবাদী} (sôtyôbādī)}{সত্যবাদী (sôtyôbādī)}]{\bn{সত্যবাদী} (sôtyôbādī) — truthful; honest in speech}

\label{lesson:BN-C350-sotyobadi}



\begin{quote}
Before the new one: say the Bengali for moral values, then the Bengali for tolerance.
\end{quote}

\subsection*{You'll want to know: \bn{সত্যবাদী}}
\textbf{\bn{সত্যবাদী}} — \emph{sôtyôbādī} — \textquotedblleft{}truthful; honest in speech\textquotedblright{}.

\begin{quote}
\bn{মীরা} \bn{একজন} \bn{সত্যবাদী} \bn{মানুষ।}
\end{quote}

\begin{quote}
\emph{mīrā ekjôn sôtyôbādī mānush.}
\end{quote}

\begin{quote}
Mira is a truthful person.
\end{quote}

\begin{cousinweb}
\textbf{\bn{সত্যবাদী}} (\emph{sôtyôbādī}) is Sanskrit \emph{satyavādin}, \emph{satya} 'truth' with \emph{vādin} 'speaking' (Monier-Williams): one whose word can be trusted.
\end{cousinweb}

\begin{grammarlens}[title={What you've built}]
One more word for giving reasons and opinions.
\end{grammarlens}

\subsection*{Your turn}
\begin{itemize}
\raggedright
  \item \emph{Say it:} \emph{sôtyôbādī}
  \item \emph{Say it:} \emph{sôtyôbādī}, once more
  \item \emph{Say it:} \emph{mīrā ekjôn sôtyôbādī mānush.} — Mira is a truthful person.
  \item \emph{Recall:} say the Bengali for moral values, then the Bengali for tolerance, then say \emph{sôtyôbādī} again
\end{itemize}

\begin{tcolorbox}[breakable,colback=teal!4,colframe=teal!35!black,title={Before you move on}]
What does \bn{সত্যবাদী} mean? (\textquotedblleft{}Truthful; honest in speech\textquotedblright{}.) Say it once more.
\end{tcolorbox}
\section[\texorpdfstring{\bn{ন্যায্য} (næjyô)}{ন্যায্য (næjyô)}]{\bn{ন্যায্য} (næjyô) — fair; rightful}

\label{lesson:BN-C350-nejyo}



\begin{quote}
Before the new one: say the Bengali for tolerance, then the Bengali for truthful.
\end{quote}

\subsection*{You'll want to know: \bn{ন্যায্য}}
\textbf{\bn{ন্যায্য}} — \emph{næjyô} — \textquotedblleft{}fair; rightful\textquotedblright{}.

\begin{quote}
\bn{আমি} \bn{ন্যায্য} \bn{দাম} \bn{চাই।}
\end{quote}

\begin{quote}
\emph{āmi næjyô dām chāi.}
\end{quote}

\begin{quote}
I want a fair price.
\end{quote}

\begin{cousinweb}
\textbf{\bn{ন্যায্য}} (\emph{næjyô}) is Sanskrit \emph{nyāyya}, right or proper (Monier-Williams), from \emph{nyāya}, what is right. A price, a demand or a share can be \bn{ন্যায্য}.
\end{cousinweb}

\begin{grammarlens}[title={What you've built}]
One more word for giving reasons and opinions.
\end{grammarlens}

\subsection*{Your turn}
\begin{itemize}
\raggedright
  \item \emph{Say it:} \emph{næjyô}
  \item \emph{Say it:} \emph{næjyô}, once more
  \item \emph{Say it:} \emph{āmi næjyô dām chāi.} — I want a fair price.
  \item \emph{Recall:} say the Bengali for tolerance, then the Bengali for truthful, then say \emph{næjyô} again
\end{itemize}

\begin{tcolorbox}[breakable,colback=teal!4,colframe=teal!35!black,title={Before you move on}]
What does \bn{ন্যায্য} mean? (\textquotedblleft{}Fair; rightful\textquotedblright{}.) Say it once more.
\end{tcolorbox}
